\chapter{Provenance, execution and immutable evidence}\label{ch:provenance}
A result is not reproducible merely because its final number is saved.
Another reader needs to know which model, protocol, implementation and
environment produced it, whether the run completed, and whether the records
were changed afterward. Provenance connects a claim to that chain.

\section{The scientific packet precedes the run}
A registered packet identifies hypotheses, domain, arms, parameters,
randomness, metrics, falsifiers and interpretation rules. It also binds the
implementation and execution environment. These decisions must precede
outcome inspection; otherwise a successful-looking configuration can be
selected from many attempts without saying so.

The present book does not fill those fields for the Gaussian programme.
It describes what must be declared. No horizon, seed family, numerical
tolerance, performance threshold or cloud resource is frozen by an example
in these chapters.

\section{Attempt, run and completion are different objects}
An invocation can fail before a scientific state transition. A campaign
can contain many registered runs within one authorized attempt. A runner
can finish producing output before the full study is durably finalized.
These distinctions prevent both inflated execution claims and unauthorized
retry disguised as routine continuation.

Historical Gate 1D-C's manifest explicitly distinguishes the runner summary
from the full-study completion sentinel \cite{gate}. That exact operational
contract belongs to that study. A future Gaussian packet must define its
own compatible mechanics and permissions, not inherit an old launch command.

\section{Content identity and scientific meaning}
Hashes identify bytes. A manifest should record the protocol, configuration,
source versions and produced artifacts so their identities can be checked.
Hashes do not certify a model's validity or a sensor's honesty. They make
it harder to substitute one object for another unnoticed.

Canonical serialization matters when equivalent-looking objects can produce
different byte sequences. Numeric representation, key ordering and schema
version are part of the identity contract. A replay should not silently
normalize a value in a way that changes the registered scientific object.

\section{What to retain in a trace}
For each event, retain enough information to reconstruct its physical
baseline, forcing, candidate/selected identities, accepted increment,
valuation inputs, owner changes, account audit and refusal state. The exact
schema is prospective, but its purpose is fixed: an independent reader must
be able to tell what happened without trusting a narrative summary alone.

Not every diagnostic belongs in the actor's view. A global evaluator may
record quantities for analysis while the selector is forbidden to read them.
Trace completeness and actor information locality are compatible when their
data paths are separated.

\section{Corrections should not erase history}
If a result artifact is wrong, the correction stage should identify the
original, the defect, the corrected interpretation and the authority for any
replacement. Rewriting the original bytes until they fit a later story
destroys the very evidence needed to understand the error.

The book revision follows the same principle. Original PDFs and historical
result artifacts are preserved. New exposition explains which model-specific
claims do not transfer, rather than editing old experimental outcomes to
look like Gaussian evidence.

\section*{Chapter recap}
Provenance makes a claim inspectable, not automatically true. Frozen
scientific choices, exact identities, completion rules and preserved failures
allow an independent review to determine what the evidence actually supports.
